% ============================================================
%  VÍAS BILIARES --- CIRUGÍAS Y PROCEDIMIENTOS
%  Licenciatura en Imagenología --- UdelaR
%  Docente: Lic. Richard González
% ============================================================

\chapter{Vías Biliares}

% ============================================================
%  PUNTOS CLAVE PARA EL EXAMEN ORAL
% ============================================================
\begin{cajakeypoints}
	\textbf{\color{azulprincipal} Puntos clave para el examen oral}
	\begin{itemize}[noitemsep, topsep=2pt]
		\item La CVL y la colangiografía abierta son procedimientos de la \textbf{vía biliar alta}; la CPRE es de la \textbf{vía biliar baja}.
		\item En la CVL el arco ingresa por el \textbf{lado izquierdo, perpendicular}; en la CPRE ingresa por el \textbf{lado derecho, perpendicular}.
		\item En la CVL el cirujano y el ayudante están del \textbf{lado izquierdo}; en la colecistectomía abierta del \textbf{lado derecho}.
		\item El catéter de Cholangiocath se introduce \textbf{2 cm} dentro del cístico y se fija con un clip de titanio.
		\item La inyección de MC en CVL: \textbf{5 cc} para la primera radiografía y \textbf{4 cc} adicionales para evaluar el vaciamiento del colédoco.
		\item En CPRE activar \textbf{ambas ``R''} al ingresar por lado derecho del paciente. La imagen se centra con la punta del endoscopio en el borde inferior, levemente a la derecha de la columna.
		\item El drenaje biliar percutáneo transhepático es paliativo para obstrucciones biliares; técnica: punción hepática $\to$ colangiografía $\to$ guía metálica $\to$ catéter de drenaje $\to$ control final.
		\item Diagnóstico diferencial fluoroscópico importante: \textbf{burbuja de aire vs.\ litiasis} en el colédoco.
	\end{itemize}
\end{cajakeypoints}

\vspace{1em}

% ============================================================
%  SECCIÓN 1 --- ANATOMÍA DE LA VÍA BILIAR
% ============================================================
\section*{Anatomía de la Vía Biliar}

\begin{cajadefinicion}
	La \textbf{vía biliar} es el conjunto de conductos que transportan la bilis desde el hígado y la vesícula biliar hasta el intestino delgado. Se divide en \textbf{intrahepática} (conductillos y conductos hepáticos derecho e izquierdo) y \textbf{extrahepática} (conducto hepático común, conducto cístico, colédoco y ampolla de Vater).
\end{cajadefinicion}

El hígado, órgano impar ubicado en el hipocondrio derecho, produce bilis con funciones metabólicas y de filtro. Posee una circulación \textbf{nutricia} (arteria hepática común) y una circulación \textbf{funcional} (sistema porta). La bilis emulsiona grasas y transporta pigmentos de desecho. El recorrido extrahepático concluye en la \textbf{ampolla hepatopancreática (de Vater)}, controlada por el esfínter de Oddi, donde confluyen el colédoco y el conducto pancreático (de Wirsung) para desembocar en el duodeno.

% ============================================================
%  SECCIÓN 2 --- CLASIFICACIÓN DE LOS PROCEDIMIENTOS
% ============================================================
\section*{Clasificación de los Procedimientos}

\subsection*{Por nivel de la vía biliar}

{\renewcommand{\arraystretch}{1.35}
	\begin{center}
		\begin{tabular}{>{\bfseries}p{3.5cm} p{4.5cm} p{6cm}}
			\hline
			Nivel & Procedimiento & Características principales \\
			\hline
			Vía biliar alta & Colangiografía laparoscópica (CVL) & Abordaje laparoscópico, catéter en cístico, MC yodado 5+4 cc. \\[0.5em]
			\hline
			Vía biliar alta & Colangiografía abierta & Misma técnica radiológica que CVL; varía el abordaje quirúrgico. \\[0.5em]
			\hline
			Vía biliar baja & CPRE & Intervención mixta endoscópica-radiológica; acceso retrógrado por papila duodenal. \\[0.5em]
			\hline
			Vía biliar intra\-hepática & Drenaje biliar percutáneo & Punción directa hepática bajo guía fluoroscópica; tratamiento paliativo de obstrucciones. \\
			\hline
		\end{tabular}
	\end{center}
}
% ============================================================
%  SECCIÓN 3 --- CVL: COLANGIOGRAFÍA LAPAROSCÓPICA
% ============================================================
\section*{Colangiografía Laparoscópica (CVL)}

\subsection*{Objetivo e indicaciones}

La CVL se realiza a criterio del cirujano durante la colecistectomía laparoscópica con cuatro objetivos principales: reconocer la \textbf{morfología de las vías biliares}, confirmar la existencia de \textbf{litiasis residual}, verificar el \textbf{pasaje del colédoco al duodeno} y detectar posibles \textbf{injurias quirúrgicas}.

\subsection*{Disposición de la sala}

{\renewcommand{\arraystretch}{1.25}
	\begin{center}
		\begin{tabular}{>{\bfseries}p{4cm} p{9.5cm}}
			\hline
			Elemento & Posición \\
			\hline
			Paciente & Decúbito dorsal, sobre mesa radiolúcida. \\[0.3em]
			\hline
			Cirujano y Ayudante 1 & Lado \textbf{izquierdo} del paciente. \\[0.3em]
			\hline
			Ayudante 2 & Lado derecho del paciente. \\[0.3em]
			\hline
			Rack de laparoscopía & Lado derecho del paciente. \\[0.3em]
			\hline
			Instrumentista y mesa & A los pies del paciente. \\[0.3em]
			\hline
			Monitor de RX & A los pies del paciente. \\[0.3em]
			\hline
			Arco en C (Licenciado) & Ingresa por el \textbf{lado izquierdo}, perpendicular. \\
			\hline
		\end{tabular}
	\end{center}
}

\subsection*{Técnica radiológica --- secuencia de pasos}

\begin{enumerate}[noitemsep]
	\item Con el abdomen insuflado, se introduce un \textbf{Jelco \#14} por debajo del reborde costal derecho a nivel de la línea medioclavicular.
	\item A través del Jelco se introduce un \textbf{catéter de Cholangiocath \#16 o \#19} (según diámetro del cístico) hacia la cavidad abdominal.
	\item Con dos pinzas de Maryland (portales \#2 y \#3) se dirige el catéter hasta el cístico.
	\item Mediante tijera recta puntiforme (portal \#2) se realiza un \textbf{pequeño corte transversal en el cístico} previamente disecado 1 cm.
	\item Se introduce el catéter \textbf{2 cm} dentro del cístico y se fija con una pinza de Maryland; el ayudante coloca un \textbf{clip de titanio} mientras inyecta solución salina para verificar paso sin fuga.
	\item Se retira el trocar \#2 y el resto del instrumental para evitar interferencias en la imagen.
	\item \textbf{Entrada del licenciado:} se posiciona el arco, se cubre y se realiza centrado previo (reperes anatómicos y materiales).
	\item \textbf{Primer disparo sin contraste} para orientar la imagen.
	\item Se coordina la inyección de \textbf{5 cc de MC yodado hidrosoluble} y se obtiene la primera radiografía.
	\item Con \textbf{4 cc adicionales} se completa el estudio para evaluar el vaciamiento del colédoco.
	\item Se guardan y documentan las imágenes.
\end{enumerate}

\vspace{0.5em}

\begin{cajakeypoints}
	\textbf{Concepto clave --- burbuja de aire vs.\ litiasis:} Una burbuja de aire aparece como imagen radiolúcida redondeada que se desplaza con los cambios de posición; la litiasis es una imagen de defecto de relleno fija. Ante la duda, cambiar la posición del paciente.
\end{cajakeypoints}

% ============================================================
%  SECCIÓN 4 --- COLECISTECTOMÍA ABIERTA
% ============================================================
\section*{Colecistectomía Abierta --- Colangiografía Intraoperatoria}

La técnica radiológica es \textbf{idéntica a la CVL}. La diferencia radica en el abordaje quirúrgico (laparotomía subcostal derecha) y en la ubicación del equipo:

\renewcommand{\arraystretch}{1.3}
\begin{center}
	\begin{tabular}{>{\bfseries}p{4cm} p{9.5cm}}
		\hline
		Elemento & Posición \\
		\hline
		Paciente & Decúbito dorsal. \\
		\hline
		Cirujano y Ayudante 1 & Lado \textbf{derecho} del paciente. \\
		\hline
		Instrumentista y mesa & Detrás del cirujano. \\
		\hline
		Ayudante 2 & Lado izquierdo del paciente. \\
		\hline
		Monitor de RX & A los pies del paciente. \\
		\hline
		Arco en C (Licenciado) & Ingresa por el \textbf{lado contrario al cirujano} (izquierdo), perpendicular. \\
		\hline
	\end{tabular}
\end{center}

\clearpage
% ============================================================
%  SECCIÓN 5 --- CPRE
% ============================================================
\section*{Colangiopancreatografía Endoscópica Retrógrada (CPRE)}

\begin{cajadefinicion}
	La \textbf{CPRE} es una intervención mixta endoscópica y radiológica, utilizada para estudiar y principalmente tratar las enfermedades de los conductos biliares y del páncreas, con acceso retrógrado a través de la papila duodenal mayor.
\end{cajadefinicion}

\subsection*{Disposición de la sala}

{\renewcommand{\arraystretch}{1.25}
	\begin{center}
		\begin{tabular}{>{\bfseries}p{4cm} p{9.5cm}}
			\hline
			Elemento & Posición \\
			\hline
			Paciente & Decúbito dorsal lateralizado a izquierda (oblicua); DLI como alternativa. \\[0.3em]
			\hline
			Cirujano, Ayudante y Mesa de instrumental & Lado \textbf{izquierdo} del paciente. \\[0.3em]
			\hline
			Monitor de endoscopía & Cabecera, lado izquierdo. \\[0.3em]
			\hline
			Monitor de RX & Detrás o al lado del monitor de endoscopía. \\[0.3em]
			\hline
			Arco en C (Licenciado) & Ingresa por el \textbf{lado derecho} del paciente, perpendicular. \\
			\hline
		\end{tabular}
	\end{center}
}

\subsection*{Técnica --- secuencia de pasos}

\begin{enumerate}[noitemsep]
	\item Anestesia superficial. Insuflación del tracto digestivo con el endoscopio.
	\item Introducción del \textbf{duodenoscopio} por vía oral: orofaringe $\to$ esófago $\to$ estómago $\to$ duodeno hasta \textbf{papila duodenal mayor}.
	\item Introducción de guía metálica en la papila: \textbf{inicio de control RX continuo}.
	\item \textbf{Cateterismo coledociano}: control RX intermitente.
	\item \textbf{Colangiografía}: coordinar inyección de MC yodado hidrosoluble --- control continuo.
	\item Según la patología: \textbf{papilotomía}, extracción de litiasis (balón con MC o canasta Dormia), colocación de stent plástico o metálico, o solo diagnóstico.
	\item \textbf{Control final}: si se colocó prótesis, disparo sin endoscopio y documentación.
\end{enumerate}

\subsection*{Orientación de la imagen en CPRE}

\begin{itemize}[noitemsep]
	\item Asegurar posición de mesa en eje Z y altura correcta.
	\item Activar \textbf{ambas ``R''} al ingresar por el lado derecho del paciente (excepción: CVL que se convierte en CPRE --- maniobra \textit{Rendez-Vous}).
	\item Paciente en \textbf{supino}: endoscopio en forma de S invertida; centrar imagen con la punta del endoscopio en el borde inferior, levemente a la derecha de la columna.
	\item Paciente en \textbf{DLI}: endoscopio en forma cóncava en su sector distal, ``mirando'' a la columna; mismo centrado que en supino.
	\item La CPRE es dinámica: no siempre hay tiempo para acomodar la imagen.
\end{itemize}

\subsection*{Intervenciones terapéuticas posibles en CPRE}

{\renewcommand{\arraystretch}{1.3}
	\begin{center}
		\begin{tabular}{>{\bfseries}p{3.5cm} p{10cm}}
			\hline
			Intervención & Descripción \\
			\hline
			Extracción de litiasis & Papilotomía + balón con MC para barrer el conducto, o canasta Dormia para capturar los cálculos. \\[0.5em]
			\hline
			Colocación de stent plástico & Tratamiento temporal de estenosis o fístulas. \\[0.5em]
			\hline
			Colocación de stent metálico autoexpansible & Tratamiento paliativo permanente en tumores irresecables. \\[0.5em]
			\hline
			Solo diagnóstico & Colangiografía sin intervención terapéutica asociada. \\
			\hline
		\end{tabular}
	\end{center}
}
% ============================================================
%  SECCIÓN 6 --- DRENAJE BILIAR PERCUTÁNEO TRANSHEPÁTICO
% ============================================================
\section*{Drenaje Biliar Percutáneo Transhepático}

\begin{cajadefinicion}
	El \textbf{drenaje biliar percutáneo transhepático} es un procedimiento terapéutico que incluye la canalización de la vía biliar intrahepática por punción directa, seguida de la utilización de guías metálicas, dilatadores y catéteres, con el objetivo de lograr el \textbf{drenaje continuo de bilis al exterior}. Es un tratamiento paliativo para patologías biliares obstructivas.
\end{cajadefinicion}

Se realiza tanto en sala de bienestar quirúrgico (BQ) como en intervencionismo. Desde el punto de vista radiológico es un procedimiento largo pero sin complicaciones habituales.

\subsection*{Disposición de la sala}

\begin{center}
	\begin{tabular}{>{\bfseries}p{4cm} p{9.5cm}}
		\toprule
		Elemento & Posición \\
		\midrule
		Paciente & Decúbito dorsal. \\[0.3em]
		Cirujano y Ayudante 1 & Lado \textbf{derecho} del paciente. \\[0.3em]
		Instrumentista y mesa & Detrás del cirujano. \\[0.3em]
		Monitor de RX & A los pies o cabeza del paciente. \\[0.3em]
		Arco en C (Licenciado) & Ingresa por el \textbf{lado izquierdo}, perpendicular. \\
		\bottomrule
	\end{tabular}
\end{center}

\subsection*{Técnica --- secuencia de pasos}

\begin{enumerate}[noitemsep]
	\item \textbf{Punción hepática} con aguja fina aspirando hasta obtener bilis (confirma canalización de la vía biliar intrahepática).
	\item \textbf{Colangiografía diagnóstica}: inyección de MC para mapear la vía biliar y localizar el nivel de obstrucción.
	\item \textbf{Colocación de guía metálica} a través de la aguja de punción.
	\item \textbf{Dilatación del trayecto} transhepático con dilatadores progresivos.
	\item \textbf{Colocación del catéter de drenaje}: puede ser externo (termina en vía biliar), interno-externo (atraviesa la obstrucción hasta el duodeno) o con stent metálico autoexpansible.
	\item \textbf{Colangiografía de control final}: verifica posición del catéter y restablecimiento del flujo biliar.
\end{enumerate}

\subsection*{Tipos de drenaje}

{\renewcommand{\arraystretch}{1.3}
	\begin{center}
		\begin{tabular}{>{\bfseries}p{3.5cm} p{5cm} p{5cm}}
			\hline
			Tipo & Descripción & Indicación \\
			\hline
			Externo & El catéter termina en la vía biliar; la bilis drena a una bolsa externa. & Primera etapa, pacientes inestables. \\[0.5em]
			\hline
			Interno-externo & El catéter atraviesa la obstrucción y llega hasta el duodeno; permite drenaje en ambas direcciones. & Pacientes con mejor condición clínica. \\[0.5em]
			\hline
			Stent metálico + posdilatación & Endoprótesis autoexpansible liberada en el sitio de estenosis; control fluoroscópico del despliegue. & Obstrucción tumoral irresecable. \\
			\hline
		\end{tabular}
	\end{center}
}

\vspace{0.5em}

\begin{cajariesgo}
	\textbf{Importante:} El calibre de la vía biliar intrahepática disminuye de forma significativa inmediatamente después de colocar el catéter de drenaje. Si no se observa esta reducción en la colangiografía de control, se debe sospechar que el drenaje no es efectivo.
\end{cajariesgo}

% ============================================================
%  SECCIÓN 7 --- COMPLICACIONES Y PROBLEMAS FRECUENTES (CVL/CPRE)
% ============================================================
\section*{Complicaciones y Problemas Frecuentes}

\subsection*{Problemas radiológicos en CVL}

\begin{itemize}[noitemsep]
	\item Al inyectar MC \textbf{no se visualiza contraste}: posible mala posición del catéter, obstrucción o fuga.
	\item \textbf{Altura de la mesa, posición del paciente y brazos} del mismo generan superposiciones o recortes de imagen.
	\item Se visualiza \textbf{MC fuera de la vía biliar}: extravasación por perforación o catéter mal posicionado.
	\item \textbf{Superposición con la columna vertebral}: ajustar angulación del arco o posición del paciente.
	\item \textbf{Calidad de imagen insuficiente}: revisar parámetros de exposición y altura del arco.
\end{itemize}

\vspace{0.5em}

\begin{cajakeypoints}
	\textbf{Diagnóstico diferencial:} La \textbf{burbuja de aire} es una imagen radiolúcida redondeada que se desplaza al cambiar la posición del paciente. La \textbf{litiasis} es un defecto de relleno fijo, con halo radiolúcido si hay contraste alrededor. Ante la duda: cambiar la posición y repetir el disparo.
\end{cajakeypoints}

% ============================================================
%  CIERRE --- ESQUEMA MENTAL FINAL
% ============================================================
\section*{Esquema Mental --- Repaso Final}

\begin{cajakeypoints}
	\begin{center}
		\textbf{\color{azulprincipal}
			Vía biliar alta (CVL/Abierta) $\rightarrow$ Vía biliar baja (CPRE) $\rightarrow$ Vía biliar intrahepática (Drenaje percutáneo)}\\[0.4em]
		\textbf{\color{naranjadestacado}
			El arco siempre ingresa perpendicular y por el lado contrario al cirujano (excepto CPRE: derecha del paciente).}\\[0.3em]
		CVL: Cholangiocath 2 cm en cístico \quad MC: 5 cc + 4 cc \quad Catéter \#16 o \#19 \\[0.3em]
		\textbf{CVL:} Arco izq. $\to$ disparo sin MC $\to$ 5 cc $\to$ 4 cc $\to$ guardar / documentar \\[0.3em]
		\textbf{CPRE:} Arco der. $\to$ ambas ``R'' $\to$ guía en papila $\to$ cateterismo $\to$ colangiografía $\to$ terapéutica $\to$ control final \\[0.3em]
		\textbf{Drenaje:} Punción $\to$ colangiografía $\to$ guía metálica $\to$ dilatación $\to$ catéter $\to$ control final
	\end{center}
\end{cajakeypoints}

\vspace{1em}
\hrule
\vspace{0.5em}
{\footnotesize \textit{Fuente: material de clase --- Lic.\ Richard González,
		Imagenología Especializada I, UdelaR 2026.}}